%==============================================================================
% documentation.tex — User Manual for the infufrgs LaTeX class
%
% Copyright (C) 2001-2026 UFRGS TeX Users Group
%
% This work may be distributed and/or modified under the conditions of
% the LaTeX Project Public License, either version 1.3c of this license
% or (at your option) any later version.
% The latest version of this license is in
%   https://www.latex-project.org/lppl.txt
% and version 1.3c or later is part of all distributions of LaTeX
% version 2005/12/01 or later.
%
% Compile with:
%   pdflatex infufrgs-documentation
%   pdflatex infufrgs-documentation
%==============================================================================
\documentclass[11pt,a4paper]{article}

\usepackage[T1]{fontenc}
\usepackage[utf8]{inputenc}
\usepackage[english]{babel}
\usepackage{geometry}
\usepackage{microtype}
\usepackage{booktabs}
\usepackage{xcolor}
\usepackage{fancyvrb}
\usepackage{array}
\usepackage{parskip}
\usepackage[colorlinks=true,linkcolor=black,urlcolor=blue,citecolor=black,
            bookmarksnumbered=true]{hyperref}

\geometry{top=25mm,bottom=25mm,left=30mm,right=25mm}

%------------------------------------------------------------------------------
% Convenience macros
%------------------------------------------------------------------------------
\newcommand{\pkg}[1]{\textsf{#1}}
\newcommand{\cls}[1]{\textsf{#1}}
\newcommand{\opt}[1]{\texttt{#1}}
\newcommand{\cmd}[1]{\texttt{\textbackslash #1}}
\newcommand{\env}[1]{\texttt{#1}}
\newcommand{\file}[1]{\texttt{#1}}
\newcommand{\meta}[1]{\textit{\textlangle#1\textrangle}}

% Inline code box
\definecolor{codebg}{rgb}{0.96,0.96,0.96}
\newcommand{\code}[1]{%
  \colorbox{codebg}{\texttt{\small #1}}}

% Verbatim environment for code listings
\DefineVerbatimEnvironment{latex}{Verbatim}{%
  frame=single,
  framesep=2mm,
  fontsize=\small,
  rulecolor=\color{gray!50},
  xleftmargin=4mm,
  xrightmargin=4mm,
  gobble=0}

% Note box
\newenvironment{note}{%
  \begin{center}
  \begin{minipage}{0.93\linewidth}
  \colorbox{yellow!20}{\parbox{\linewidth-2\fboxsep}{%
  \smallskip\textbf{Note:}\enspace\ignorespaces
}{%
  \smallskip}}
  \end{minipage}
  \end{center}}

%------------------------------------------------------------------------------
% Document metadata
%------------------------------------------------------------------------------
\title{%
  \pkg{infufrgs} — LaTeX Class for\\[4pt]
  INF/UFRGS Academic Documents\\[8pt]
  \large\normalfont User Manual\\[4pt]
  \large\normalfont Version~5.0.0 (2026-03-15)}
\author{UFRGS TeX Users Group\\
  \small Institute of Informatics — UFRGS, Porto Alegre, Brazil\\
  \small\url{https://www.inf.ufrgs.br/utug}}
\date{}

%==============================================================================
\begin{document}
%==============================================================================

\maketitle
\thispagestyle{empty}

\begin{abstract}
  \noindent
  The \cls{infufrgs} document class typesets theses, dissertations,
  monographs, research reports, and related academic documents following
  the typographic and bibliographic standards required by the Institute
  of Informatics (INF) of the Universidade Federal do Rio Grande do Sul
  (UFRGS), Porto Alegre, Brazil.  This manual describes version~5.0.0,
  a self-contained single-file consolidation released under the
  LaTeX Project Public Licence~1.3c.
\end{abstract}

\tableofcontents

%==============================================================================
\section{Introduction}
%==============================================================================

The \cls{infufrgs} class has been maintained by the UFRGS TeX Users Group
since~2001.  It builds on the standard \cls{report} class and provides:
%
\begin{itemize}
  \item automatic typesetting of the title page and CIP
        (Catalogação na Publicação) block;
  \item front-matter environments for abstracts in two languages,
        dedications, acknowledgements, and lists of abbreviations
        and symbols;
  \item chapter, section, and caption formatting that complies with
        the ABNT standards adopted by INF/UFRGS;
  \item built-in support for ABNT-style bibliographic citations via
        the \pkg{abntex2cite} package;
  \item support for several undergraduate and graduate programmes
        at INF/UFRGS and the Microelectronics and Letters
        Post-Graduate Programmes;
  \item optional theorem/definition environments (\opt{formais} option);
  \item support for UNIBRAL (UFRGS–TU\,Berlin) double-degree documents
        (\opt{unibral} option).
\end{itemize}
%
Version~5.0.0 is a major modernisation: all auxiliary files
(\file{iidefs.sty}, \file{*.def}, \file{*.babel}) have been
consolidated into the single file \file{infufrgs.cls}.

%==============================================================================
\section{Requirements}
\label{sec:requirements}
%==============================================================================

\cls{infufrgs} requires a reasonably modern TeX installation.
TeX~Live~2021 or later (or MiK\TeX{} with all packages updated) is
recommended.  The following packages must be available; all of them
are part of standard TeX~Live and MiK\TeX{} distributions.

\begin{center}
\begin{tabular}{@{}ll@{}}
  \toprule
  Package & Purpose \\
  \midrule
  \pkg{xspace}       & Smart spacing after macros \\
  \pkg{keyval}       & Key–value option parsing \\
  \pkg{iftex}        & Engine detection (\cmd{ifpdftex}, \cmd{ifxetex}, \cmd{ifluatex}) \\
  \pkg{babel}        & Language support (Brazilian Portuguese and English) \\
  \pkg{lastpage}     & \verb|\pageref{LastPage}| for CIP \\
  \pkg{indentfirst}  & Indent first paragraph of each section \\
  \pkg{setspace}     & Line-spacing control \\
  \pkg{titlesec}     & Section heading formatting \\
  \pkg{tocloft}      & Table-of-contents formatting \\
  \pkg{fancyhdr}     & Page header/footer \\
  \pkg{xcolor}       & Colour support \\
  \pkg{graphicx}     & Figure inclusion \\
  \pkg{caption}      & Caption formatting \\
  \pkg{hyperref}     & PDF hyperlinks and metadata \\
  \pkg{hypcap}       & Correct hyperlink anchors for floats \\
  \pkg{geometry}     & Page layout \\
  \pkg{newtxtext}, \pkg{newtxmath} & Times-based fonts (pdf\LaTeX{} only) \\
  \pkg{fontspec}     & OpenType fonts (Xe\LaTeX/Lua\LaTeX{} only) \\
  \pkg{abntex2cite}  & ABNT citation style (part of \pkg{abntex2}) \\
  \bottomrule
\end{tabular}
\end{center}

The \pkg{abntex2} bundle is shipped with TeX~Live and MiK\TeX{}.
Font packages \pkg{newtxtext}/\pkg{newtxmath} are available in TeX~Live
as \texttt{tex-newtx}.

%==============================================================================
\section{Installation}
\label{sec:installation}
%==============================================================================

Copy \file{infufrgs.cls} to a directory where \TeX{} can find it.
Common choices are:
%
\begin{itemize}
  \item the same directory as your document (simplest);
  \item \texttt{TEXMFHOME/tex/latex/infufrgs/} for a personal
        installation (run \texttt{texhash} afterwards).
\end{itemize}
%
When installed from CTAN via TeX~Live (\texttt{tlmgr install infufrgs})
or MiK\TeX, no manual step is required.

%==============================================================================
\section{Class Options}
\label{sec:options}
%==============================================================================

Options are supplied as a comma-separated list in the
\cmd{documentclass} argument:

\begin{latex}
\documentclass[<programme>,<document-type>,...]{infufrgs}
\end{latex}

\subsection{Programme options}
\label{sec:progopts}

Exactly one programme option should be given.

\begin{center}
\begin{tabular}{@{}lp{0.70\linewidth}@{}}
  \toprule
  Option      & Programme \\
  \midrule
  \opt{cic}     & Ciência da Computação (undergraduate) \\
  \opt{ecp}     & Engenharia de Computação (undergraduate) \\
  \opt{ppgc}    & Programa de Pós-Graduação em Computação \\
  \opt{pgmicro} & Programa de Pós-Graduação em Microeletrônica \\
  \opt{ppgl}    & Programa de Pós-Graduação em Letras \\
  \bottomrule
\end{tabular}
\end{center}

\subsection{Document-type options}
\label{sec:doctypeopts}

Exactly one document-type option should be given.

\begin{center}
\begin{tabular}{@{}lp{0.68\linewidth}@{}}
  \toprule
  Option(s)              & Document type \\
  \midrule
  \opt{tc} / \opt{dipl}             & Trabalho de Conclusão (bachelor's monograph) \\
  \opt{diss} / \opt{mestrado}       & Dissertação de Mestrado (Master's thesis) \\
  \opt{tese} / \opt{doutorado}      & Tese de Doutorado (Doctoral thesis) \\
  \opt{espec}                       & Monografia de Especialização \\
  \opt{ti}                          & Trabalho Individual \\
  \opt{rp}                          & Relatório de Pesquisa \\
  \opt{pep}                         & Plano de Estudos e Pesquisa \\
  \opt{prop-tese}                   & Proposta de Tese de Doutorado \\
  \opt{plano-doutorado}             & Plano de Doutorado \\
  \bottomrule
\end{tabular}
\end{center}

\subsection{Other options}
\label{sec:otheropts}

\begin{center}
\begin{tabular}{@{}lp{0.68\linewidth}@{}}
  \toprule
  Option              & Effect \\
  \midrule
  \opt{english}       & Typeset the document in English (headings, captions, etc.) \\
  \opt{oneside}       & Single-sided layout (default is two-sided) \\
  \opt{noabntcite}    & Do not automatically load \pkg{abntex2cite};
                        load it yourself or use a different citation package \\
  \opt{nominatalocal} & Read nominata data from \file{nominatalocal.def}
                        instead of the built-in defaults
                        (see Section~\ref{sec:nominata}) \\
  \opt{formais}       & Load theorem/definition environments
                        (see Section~\ref{sec:formais}) \\
  \opt{unibral}       & Enable UNIBRAL double-degree support
                        (see Section~\ref{sec:unibral}) \\
  \bottomrule
\end{tabular}
\end{center}

Any option not recognised by \cls{infufrgs} is passed through to the
underlying \cls{report} class.

%==============================================================================
\section{Document Preamble — Metadata Commands}
\label{sec:preamble}
%==============================================================================

All the following commands should appear in the preamble, i.e., between
\cmd{documentclass} and \cmd{begin\{document\}}.

\subsection{Title and translated title}

\begin{latex}
\title{<title of the document>}
\translatedtitle{<title in the other language>}
\end{latex}

\cmd{title} sets the document title.  \cmd{translatedtitle} provides
the title in the second language (Portuguese if the document is in
English, English if the document is in Portuguese).  The translated
title is displayed on the \env{translatedabstract} page.

\subsection{Author(s)}

\begin{latex}
\author{<surname>}{<first name(s)>}
\end{latex}

Up to four authors are supported; call \cmd{author} once for each.
The arguments are kept separate so that the CIP block (Catalogação na
Publicação) can format the name correctly.  Example:

\begin{latex}
\author{Flaumann}{Fritz Gutenberg}
\author{Flaumann}{Frida Gutenberg}
\end{latex}

\subsection{Advisor and co-advisor}

\begin{latex}
\advisor[<title>]{<surname>}{<first name(s)>}
\coadvisor[<title>]{<surname>}{<first name(s)>}
\end{latex}

The optional \meta{title} argument is printed before the name
(e.g., \texttt{Prof.\textasciitilde Dr.}).  \cmd{coadvisor} is optional.

\subsection{Date}

\begin{latex}
\date{<month>}{<year>}
\end{latex}

If omitted, the current month and year are used.
For a Portuguese document use the month name in Portuguese
(e.g., \texttt{maio}); for an English document, in English
(e.g., \texttt{May}).

\subsection{Location}

\begin{latex}
\location{<city>}{<state abbreviation>}
\end{latex}

Defaults to \texttt{Porto\textasciitilde Alegre} / \texttt{RS}.

\subsection{Keywords}

\begin{latex}
\keyword{<keyword in the primary language>}
\translatedkeyword{<keyword in the secondary language>}
\end{latex}

Call these commands once per keyword (up to ten each).  Keywords appear
after the abstract text and in the PDF metadata.

\subsection{Department (optional)}

\begin{latex}
\dept{<department code>}
\end{latex}

Currently recognised codes: \opt{INA} (Departamento de Informática
Aplicada) and \opt{INT} (Departamento de Informática Teórica).

\subsection{Additional commands for specific document types}

\begin{center}
\begin{tabular}{@{}lll@{}}
  \toprule
  Document type & Command & Argument \\
  \midrule
  \opt{ti}    & \cmd{ti}\texttt{[release]\{number\}}
              & Release (default \texttt{I}) and TI number \\
  \opt{rp}    & \cmd{rp}\texttt{\{number\}}
              & RP number \\
  \opt{rp}    & \cmd{financ}\texttt{\{agencies\}}
              & Funding agencies (optional) \\
  \opt{espec} & \cmd{espec}\texttt{\{specialisation name\}}
              & Name of the specialisation programme \\
  \opt{espec} & \cmd{coord}\texttt{[title]\{surname\}\{first name\}}
              & Course coordinator \\
  \bottomrule
\end{tabular}
\end{center}

%==============================================================================
\section{Front Matter}
\label{sec:frontmatter}
%==============================================================================

\subsection{Title page and CIP/nominata}

\begin{latex}
\begin{document}
\maketitle
\end{latex}

\cmd{maketitle} generates:
\begin{enumerate}
  \item the title page, and
  \item for Master/Doctoral/Specialisation documents: a CIP
        (Catalogação na Publicação) block followed by the
        \emph{nominata} (administrative officers listing);
        for other document types: only the nominata.
\end{enumerate}

\subsection{Dedication}

\begin{latex}
\begin{dedicatoria}
  To my family.
\end{dedicatoria}
\end{latex}

Produces a page with the dedication text flush-right at the bottom.
No section heading is added.

\subsection{Acknowledgements}

\begin{latex}
\begin{agradecimentos}
  I thank everyone who contributed\ldots
\end{agradecimentos}
\end{latex}

Produces an unnumbered chapter titled \textit{Agradecimentos}.

%==============================================================================
\section{Abstracts}
\label{sec:abstracts}
%==============================================================================

\subsection{Primary-language abstract}

\begin{latex}
\begin{abstract}
  <Abstract text.>
\end{abstract}
\end{latex}

Produces an unnumbered chapter titled \textit{Resumo} (or
\textit{Abstract} for English documents).  The keywords defined with
\cmd{keyword} are appended automatically.

\subsection{Secondary-language abstract}

\begin{latex}
\begin{translatedabstract}
  <Abstract in the other language.>
\end{translatedabstract}
\end{latex}

The environment automatically switches to the other language
(English--Brazilian Portuguese).  The title of the document in the
other language, set with \cmd{translatedtitle}, is displayed above
the abstract text.  The keywords defined with \cmd{translatedkeyword}
are appended automatically.

\bigskip
\noindent\textbf{Deprecated:} the \env{englishabstract} environment
from earlier versions is still accepted for backwards compatibility
but produces a warning.  Replace it with \env{translatedabstract}.

%==============================================================================
\section{Lists}
\label{sec:lists}
%==============================================================================

The standard list commands work as usual:

\begin{latex}
\tableofcontents
\listoffigures
\listoftables
\end{latex}

\subsection{List of abbreviations and acronyms}

\begin{latex}
\begin{listofabbrv}{<widest-label>}
  \item[ABNT]  Associação Brasileira de Normas Técnicas
  \item[INF]   Instituto de Informática
  \item[UFRGS] Universidade Federal do Rio Grande do Sul
\end{listofabbrv}
\end{latex}

The mandatory argument \meta{widest-label} (e.g., \texttt{UFRGS}) is used
to compute the column width so that all abbreviations align.

\subsection{List of symbols}

\begin{latex}
\begin{listofsymbols}{<widest-symbol>}
  \item[$\alpha$] Alpha coefficient
  \item[$\beta$]  Beta coefficient
\end{listofsymbols}
\end{latex}

Same structure as \env{listofabbrv}.

%==============================================================================
\section{Main Text}
\label{sec:maintext}
%==============================================================================

\subsection{Sectioning}

The class provides all standard \cls{report} sectioning commands:
\cmd{chapter}, \cmd{section}, \cmd{subsection}, \cmd{subsubsection},
and \cmd{paragraph}.  Section depth is set to~5 (\cmd{secnumdepth}).

Chapter headings are typeset in bold, 12\,pt, uppercase, flush-left.
Sections and subsections are bold; subsubsections are italic.

Chapters and sections are recorded in the table of contents with
uppercase names for chapters.

\subsection{Figures and tables}

Standard \env{figure} and \env{table} float environments are used.
Captions are formatted with an em-dash separator, centred,
in 10\,pt type.

\begin{latex}
\begin{figure}[htbp]
  \centering
  \includegraphics[width=0.7\textwidth]{filename}
  \caption{Caption text.}
  \label{fig:label}
\end{figure}
\end{latex}

The \cmd{legend} command produces an unnumbered caption:

\begin{latex}
\legend{Source: <reference>}
\end{latex}

\subsection{Long quotations}

The \env{quote} environment is redefined to comply with ABNT formatting:
4\,cm left indentation, 10\,pt font, single spacing.

\begin{latex}
\begin{quote}
  Long quotation text of more than three lines goes here.
\end{quote}
\end{latex}

\subsection{Lists}

The \env{enumerate}, \env{itemize}, and \env{description} environments
are redefined with reduced vertical spacing to match the ABNT style.

%==============================================================================
\section{Back Matter}
\label{sec:backmatter}
%==============================================================================

\subsection{Appendices}

\begin{latex}
\appendix
\chapter{First Appendix}
...
\end{latex}

\cmd{appendix} resets the chapter counter and switches to alphabetic
numbering.  The optional argument overrides the heading label
(default: \textit{Apêndice} / \textit{Appendix}).

\subsection{Annexes}

\begin{latex}
\annex
\chapter{First Annex}
...
\end{latex}

\cmd{annex} is similar to \cmd{appendix} but uses the heading label
\textit{Anexo} / \textit{Annex}.  In Brazilian academic convention,
appendices contain material produced by the author; annexes contain
third-party material reproduced for documentation.

\subsection{Bibliography}

The class automatically loads \pkg{abntex2cite} (unless the
\opt{noabntcite} option is given).  Use one of the ABNT styles:

\begin{latex}
\bibliographystyle{abntex2-alf}   % author-year citations
%\bibliographystyle{abntex2-num}  % numeric citations
\bibliography{<bib file(s)>}
\end{latex}

%==============================================================================
\section{Special Features}
%==============================================================================

\subsection{Formal environments (\opt{formais} option)}
\label{sec:formais}

Loading \opt{formais} adds a collection of theorem-like environments
and two additional lists.

\subsubsection{Theorem-like environments}

\begin{center}
\begin{tabular}{@{}ll@{}}
  \toprule
  Environment & Name \\
  \midrule
  \env{theorem}      & Theorem / Teorema \\
  \env{lemma}        & Lemma / Lema \\
  \env{corollary}    & Corollary / Corolário \\
  \env{proposition}  & Proposition / Proposição \\
  \env{conjecture}   & Conjecture / Conjectura \\
  \env{definition}   & Definition / Definição \\
  \env{example}      & Example / Exemplo \\
  \env{exercise}     & Exercise / Exercício \\
  \env{property}     & Property / Propriedade \\
  \env{remark}       & Remark / Observação \\
  \bottomrule
\end{tabular}
\end{center}

All environments take an optional name argument.  Example:

\begin{latex}
\begin{theorem}[Pythagorean theorem]
  In a right triangle, $a^2 + b^2 = c^2$.
\end{theorem}
\begin{proof}
  See any introductory geometry text.
\end{proof}
\end{latex}

\subsubsection{Lists of definitions and theorems}

\begin{latex}
\listofdefinitions
\listoftheorems
\end{latex}

These produce numbered lists analogous to \cmd{listoffigures}.

\subsection{UNIBRAL double-degree (\opt{unibral} option)}
\label{sec:unibral}

The \opt{unibral} option activates support for the UNIBRAL programme
(joint double-degree between UFRGS and Technische Universität Berlin).

\subsubsection{Additional commands}

\begin{latex}
\advisorDE[<title>]{<surname>}{<first name>}
\coadvisorDE[<title>]{<surname>}{<first name>}
\end{latex}

These set the German advisor and co-advisor, which appear on the title
page alongside the Brazilian advisor(s).

\subsubsection{Extended summary}

\begin{latex}
\extendedsummary
<Portuguese extended summary for UFRGS.>
\end{latex}

Produces the \emph{Resumo Estendido} chapter required by UFRGS for
double-degree theses.

%==============================================================================
\section{Nominata Customisation}
\label{sec:nominata}
%==============================================================================

The nominata is the list of UFRGS administrative officers printed on
the CIP/nominata page.  The class ships with defaults reflecting the
current officeholders; however, these data change every few years.

\subsection{Overriding individual fields}

Any nominata field can be updated in the preamble with
\cmd{renewcommand}:

\begin{latex}
\renewcommand{\nominataReit}{Prof\textsuperscript{a}.~Ana Silva}
\renewcommand{\nominataReitname}{Reitora}
\end{latex}

The available fields are listed in the table below.

\begin{center}
\footnotesize
\begin{tabular}{@{}ll@{}}
  \toprule
  Command & Content \\
  \midrule
  \cmd{nominataReit}          & Rector \\
  \cmd{nominataReitname}      & Rector's title \\
  \cmd{nominataPRCA}          & Vice-rector \\
  \cmd{nominataPRCAname}      & Vice-rector's title \\
  \cmd{nominataPRAPG}         & Pro-rector for Graduate Studies \\
  \cmd{nominataPRAPGname}     & Pro-rector's title \\
  \cmd{nominataPRAG}          & Pro-rector for Undergraduate Studies \\
  \cmd{nominataPRAGname}      & Pro-rector's title \\
  \cmd{nominataDir}           & Director of INF \\
  \cmd{nominataDirname}       & Director's title \\
  \cmd{nominataCoordPPGC}     & PPGC coordinator \\
  \cmd{nominataCoordnamePPGC} & PPGC coordinator title \\
  \cmd{nominataCoordPGMICRO}  & PGMICRO coordinator \\
  \cmd{nominataCoordnamePGMICRO} & PGMICRO coordinator title \\
  \cmd{nominataBibchefe}      & Head librarian \\
  \cmd{nominataBibchefename}  & Head librarian's title \\
  \cmd{nominataChefeINA}      & Head of INA department \\
  \cmd{nominataChefeINAname}  & INA department head title \\
  \cmd{nominataChefeINT}      & Head of INT department \\
  \cmd{nominataChefeINTname}  & INT department head title \\
  \cmd{nominataCoordCIC}      & CIC course coordinator \\
  \cmd{nominataCoordCICname}  & CIC coordinator title \\
  \cmd{nominataCoordECP}      & ECP course coordinator \\
  \cmd{nominataCoordECPname}  & ECP coordinator title \\
  \bottomrule
\end{tabular}
\end{center}

\subsection{External file (\opt{nominatalocal} option)}

If the \opt{nominatalocal} class option is given, all nominata data
are read from the file \file{nominatalocal.def} (searched via the
normal \TeX{} input path).  This file should contain
\cmd{providecommand} or \cmd{renewcommand} definitions for the fields
listed above.  This mechanism is convenient when the same nominata
data are shared by many documents in an institutional repository.

%==============================================================================
\section{A Complete Minimal Example}
\label{sec:example}
%==============================================================================

The following skeleton compiles a Master's dissertation for PPGC:

\begin{latex}
\documentclass[ppgc,diss]{infufrgs}

\usepackage[utf8]{inputenc}   % for pdfLaTeX; not needed for XeLaTeX/LuaLaTeX
\usepackage{graphicx}

\title{Title of the Dissertation}
\translatedtitle{Título da Dissertação}
\author{Surname}{First Name}
\advisor[Prof.~Dr.]{AdvisorSurname}{AdvisorFirst}
% \coadvisor[Prof.~Dr.]{CoAdvisorSurname}{CoAdvisorFirst}
\date{March}{2026}

\keyword{Keyword one}
\keyword{Keyword two}
\translatedkeyword{Primeira palavra-chave}
\translatedkeyword{Segunda palavra-chave}

\begin{document}

\maketitle

\begin{dedicatoria}
  To my family.
\end{dedicatoria}

\begin{agradecimentos}
  I thank everyone\ldots
\end{agradecimentos}

\begin{abstract}
  Abstract text in English.
\end{abstract}

\begin{translatedabstract}
  Texto do resumo em português.
\end{translatedabstract}

\tableofcontents
\listoffigures
\listoftables

\begin{listofabbrv}{UFRGS}
  \item[ABNT]  Associação Brasileira de Normas Técnicas
  \item[UFRGS] Universidade Federal do Rio Grande do Sul
\end{listofabbrv}

\chapter{Introduction}
...

\appendix
\chapter{Supporting Material}
...

\bibliographystyle{abntex2-alf}
\bibliography{references}

\end{document}
\end{latex}

A fully compiled example is provided in the file \file{example.tex}
(together with \file{example.bib}).

%==============================================================================
\section{Compilation}
\label{sec:compilation}
%==============================================================================

\cls{infufrgs} works with pdf\LaTeX, Xe\LaTeX, and Lua\LaTeX.

For pdf\LaTeX{}, declare the input encoding in the preamble:
\begin{latex}
\usepackage[utf8]{inputenc}
\end{latex}
This is not required for Xe\LaTeX{} or Lua\LaTeX, which handle UTF-8
natively.

The typical compilation sequence for a document with a bibliography is:

\begin{latex}
pdflatex <document>
bibtex   <document>
pdflatex <document>
pdflatex <document>
\end{latex}

Or, using \texttt{latexmk}:

\begin{latex}
latexmk -pdf <document>
\end{latex}

%==============================================================================
\section{Package Files}
\label{sec:files}
%==============================================================================

The distribution contains the following files:

\begin{center}
\begin{tabular}{@{}ll@{}}
  \toprule
  File & Purpose \\
  \midrule
  \file{infufrgs.cls}    & The document class (single self-contained file) \\
  \file{infufrgs-doc.tex} & This user manual (source) \\
  \file{infufrgs-example.tex}     & Complete compilable usage example \\
  \file{infufrgs-example.bib}     & Bibliography database for the example \\
  \file{README.md}       & Short introduction and quick-start guide \\
  \bottomrule
\end{tabular}
\end{center}

%==============================================================================
\section{Licence}
\label{sec:licence}
%==============================================================================

This work may be distributed and/or modified under the conditions of the
\textbf{LaTeX Project Public License} (LPPL), either version~1.3c of
this licence or (at your option) any later version.  The latest version
of this licence is at \url{https://www.latex-project.org/lppl.txt} and
version~1.3c or later is part of all distributions of \LaTeX{} version
2005/12/01 or later.

This work has the LPPL maintenance status \textit{maintained}.

\medskip
\noindent
Copyright \textcopyright\ 2001-2026 UFRGS TeX Users Group\\
Institute of Informatics — UFRGS, Porto Alegre, Brazil\\
\url{https://www.inf.ufrgs.br/utug}

%==============================================================================
\section{Version History}
\label{sec:history}
%==============================================================================

\begin{description}
  \item[5.0.0 2026-03-15]
    Major modernisation.  All auxiliary files
    (\file{iidefs.sty}, \file{*.def}, \file{*.babel}) consolidated
    into \file{infufrgs.cls}.
    Replaced deprecated \pkg{tocstyle} with \pkg{tocloft}.
    Replaced obsolete \pkg{times} font package with
    \pkg{newtxtext}/\pkg{newtxmath} (pdf\LaTeX) and \pkg{fontspec}
    with TeX Gyre Termes (Xe/Lua\LaTeX).
    Added \pkg{iftex} for engine detection.
    Licence changed from GPL-2 to LPPL~1.3c.
    TOC uppercase chapter titles now use \cmd{texorpdfstring} so
    PDF bookmarks retain mixed-case titles.
    Added \env{translatedabstract} environment; deprecated
    \env{englishabstract} (kept for backwards compatibility).

  \item[4.x series (2010–2024)]
    Various maintenance updates: nominata updates, minor bug fixes,
    compatibility with evolving \LaTeX{} distributions.

  \item[1.0 (2001)]
    Initial release by the UFRGS TeX Users Group.
\end{description}

%==============================================================================
\end{document}
%==============================================================================
